\addcontentsline{toc}{chapter}{Introduction générale}
\chapter*{Introduction générale}
Dans un monde où les transactions électroniques sont devenues essentielles, les systèmes de paiement doivent garantir des échanges sûrs, fiables et conformes aux normes internationales. Parmi ces normes, le protocole ISO8583 est largement utilisé pour structurer et transmettre les messages de transactions financières entre terminaux de paiement et systèmes bancaires. Toutefois, tester ces échanges en environnement réel est souvent complexe, coûteux et risqué. D’où l’importance de disposer d’un outil capable de simuler, tester et valider ces transactions avant leur mise en production.

Dans ce cadre s’inscrit notre projet, réalisé durant un stage chez HPS (Hightech Payment Systems), acteur marocain de renommée internationale dans les solutions de paiement électronique, avec son produit phare PowerCARD. L’objectif est de concevoir un simulateur ISO8583 basé sur une architecture microservices, capable de tester, valider et analyser les transactions financières de manière automatisée, avec des fonctionnalités telles que la construction, la validation, le chiffrement et la communication asynchrone des messages.

Grâce à une architecture modulaire reposant sur Spring Boot, Spring Cloud, JPOS et Angular, ce simulateur vise à améliorer la fiabilité et la qualité des systèmes de paiement. Ce projet représente une avancée stratégique pour HPS, en renforçant ses capacités de support et en optimisant les processus de validation. Il m'a également permis de mettre en pratique mes compétences techniques tout en contribuant à l'amélioration de l'efficacité de l’infrastructure de test.

Ce rapport est structuré en cinq chapitres. Il commence par une présentation du contexte général du projet, en abordant les enjeux liés aux transactions électroniques et les limites des tests en environnement réel. Le deuxième chapitre présente la conduite du projet ainsi que l’analyse des besoins fonctionnels et techniques. Le troisième chapitre est consacré à l’architecture technique et à l’implémentation du simulateur. Le quatrième chapitre propose une évaluation des résultats obtenus et des pistes d’amélioration. Enfin, le cinquième chapitre conclut le rapport en résumant les principales réalisations et en proposant des perspectives d’évolution.

